\section{Формальная постановка задачи}



Пусть задан запрос, затрагивающий $n$ базовых отношений $\mathcal{R} = \{R_1, \ldots, R_n\}$, и
множество предикатов соединения $\mathcal{P}$, определяющее граф соединений
$G = (\mathcal{R},\, \mathcal{E})$, где ребро $(R_i, R_j) \in \mathcal{E}$ существует тогда и только тогда,
когда имеется предикат, связывающий $R_i$ и $R_j$.

\textbf{План исполнения} запроса~--- это план отношения для всего множества
$\mathcal{R}$: бинарное дерево $T$ с $n$ листьями, помеченными
отношениями из $\mathcal{R}$, в котором каждый внутренний узел помечен
парой $(\Join,\, op)$~--- логической операцией соединения своих двух
поддеревьев и физическим алгоритмом $op$, реализующим это соединение
(например, хеш-соединение, соединение слиянием, вложенное сканирование, вложенное сканирование с помощью индекса).
Физический алгоритм $op$ допустим в узле только при выполнении его
предусловий (например, наличие индекса на
одном из входных отношении или отсортированности входов).
Обозначим $\mathcal{T}_n$~--- множество всех допустимых планов
для данного запроса.
\par

\textbf{Алгоритм планирования.}
Алгоритм планирования~--- это отображение
$\mathcal{A}\colon (\mathcal{R},\, \mathcal{P},\, \Sigma,\ \mathcal(D)) \to \mathcal{T}_n$,
где $\Sigma$~--- множество доступных статистик (гистограммы, список самых частых значений,
оценки уникальных значений для каждого атрибута), $\mathcal(D)$ ~--- данные.
Результат $T_{\mathcal{A}} = \mathcal{A}(\mathcal{R}, \mathcal{P}, \Sigma) \in \mathcal{T}_n$~---
план исполнения.
Работа алгоритма $\mathcal{A}$ характеризуется двумя величинами:
\begin{itemize}
    \item $t_{\mathrm{plan}}(\mathcal{A})$~--- время работы алгоритма планирования (время перебора, оценки и выбора плана);
    \item $t_{\mathrm{exec}}(\mathcal{A})$~--- время исполнения выбранного плана $T_{\mathcal{A}}$ на данных.
\end{itemize}

Множество $\mathcal{T}_n$ включает как комбинаторику логических порядков
соединений (форма дерева и распределение отношений по листьям), так и
назначение физических операторов внутренним узлам. В  работе выбор физического оператора и проверка его предусловий делегируются 
процедуре построения путей PostgreSQL; предлагаемые методы действуют на пространстве \emph{логических} порядков соединений,
сокращая или упорядочивая его перед передачей в эту процедуру.
\par

Время исполнения зависит от качества выбранного плана: чем ближе $T_{\mathcal{A}}$ к истинному оптимуму
по фактическому времени, тем меньше $t_{\mathrm{exec}}$.
Время планирования определяется алгоритмом $\mathcal{A}$: точные методы (полный DP-перебор) дают
малое $t_{\mathrm{exec}}$, но большое $t_{\mathrm{plan}}$ при росте $n$;
быстрые эвристики сокращают $t_{\mathrm{plan}}$, но могут увеличить $t_{\mathrm{exec}}$.
\par

\textbf{Целевая функция.}
Суммарное время обработки запроса:
\[
    t_{\mathrm{total}}(\mathcal{A}) \;=\; t_{\mathrm{plan}}(\mathcal{A}) \;+\; t_{\mathrm{exec}}(\mathcal{A}) \;\longrightarrow\; \min_{\mathcal{A}}.
\]
Задача состоит в построении алгоритма планирования $\mathcal{A}^*$,
минимизирующего $t_{\mathrm{total}}$ на заданном классе запросов $\mathcal{Q}$:
\[
    \mathcal{A}^* \;=\; \arg\min_{\mathcal{A}} \sum_{q \in \mathcal{Q}} t_{\mathrm{total}}(\mathcal{A},\, q).
\]